\section{Кинематика}
%%%%%%%%%%%%%%%%%%%%%%%%%%%%%%%%%%%%%%%%%%%%%%%%%%%%
\subsection{Графики}
%%%%%%%%%%%%%%%%%%%%%%%%%%%%%%%%%%%%%%%%%%%%%%%%%%%%
%%%%%%%%%%%%%%%%%%%%%%%%%%%%%%%%%%%%%%%%%%%%%%%%%%%%
\z[{\begin{wrapfigure}[6]{r}{0.3\textwidth}%
		\vspace{-15pt}%
		\begin{flushright}%
 			\includegraphics[width=0.3\textwidth]{2.1.1}%
		\end{flushright}%
	\end{wrapfigure}%
	}]%
{По приведённым зависимостям координаты движущегося вдоль оси $O_x$ тела от времени $x(t)$ постройте графики зависимостей от времени для: пройденного телом пути $s(t)$, модуля скорости тела $v(t)$, проекции скорости $v_x(t)$, средней путевой скорости $v_{\text{ср}}(t)$ и зависимости от координаты $x$ для: модуля скорости тела $v(x)$ и средней путевой скорости $v_{\text{ср}}(x)$.}
%
{}

%%%%%%%%%%%%%%%%%%%%%%%%%%%%%%%%%%%%%%%%%%%%%%%%%%%%
\z[{\begin{wrapfigure}[5]{r}{0.3\textwidth}%
		\vspace{-25pt}%
		\begin{flushright}%
 			\includegraphics[width=0.3\textwidth]{2.1.2}%
		\end{flushright}%
	\end{wrapfigure}%
	}]%
{Экспериментатор Глюк исследовал немного необычные графики равномерного движения двух прямолинейно движущихся тел (см. рисунок). Определите, какое из тел (I или II) двигалось с большей скоростью. Во сколько раз отличались скорости тел?}
%
{У второго в $1.5$ раза}

%%%%%%%%%%%%%%%%%%%%%%%%%%%%%%%%%%%%%%%%%%%%%%%%%%%%
\z[{\begin{wrapfigure}[5]{r}{0.3\textwidth}%
		\vspace{-5pt}%
		\begin{flushright}%
 			\includegraphics[width=0.3\textwidth]{2.1.3}%
		\end{flushright}%
	\end{wrapfigure}%
	}]%
{Муха движется вдоль оси $O_x$. В начальный момент времени она находилась в точке с координатой $x = 5\,\text{м}$. На рисунке изображена зависимость проекции скорости $v_x$ мухи от времени $t$. Чему равна её координата в моменты времени $t_1 = 2\,\text{с}$ и $t_1 = 4\,\text{с}$?}
%
{$9\,\text{м}$ и $5\,\text{м}$}

%%%%%%%%%%%%%%%%%%%%%%%%%%%%%%%%%%%%%%%%%%%%%%%%%%%%
\z[{\begin{wrapfigure}[5]{r}{0.3\textwidth}%
		\vspace{+5pt}%
		\begin{flushright}%
 			\includegraphics[width=0.3\textwidth]{2.1.4}%
		\end{flushright}%
	\end{wrapfigure}%
	}]%
{Из двух городов навстречу с постоянными скоростями движутся два автомобиля. На графике приведена зависимость расстояния между автомобилями от времени их движения. Определите модуль скорости первого автомобиля в системе отсчёта, связанной со вторым автомобилем.}
%
{$100\,\text{км/ч}$}

%%%%%%%%%%%%%%%%%%%%%%%%%%%%%%%%%%%%%%%%%%%%%%%%%%%%
\z[{\begin{wrapfigure}[6]{r}{0.3\textwidth}%
		\vspace{-10pt}%
		\begin{flushright}%
 			\includegraphics[width=0.3\textwidth]{2.1.5}%
		\end{flushright}%
	\end{wrapfigure}%
	}]%
{Пчела отправилась из улья за мёдом. Сначала она летела медленно в поисках мёда, затем собирала его, а затем быстро летела обратно, но не по кратчайшему пути. На графике приведена зависимость скорости пчелы от времени её движения. Постройте график зависимости скорости пчелы от пройденного ей расстояния.}
%
{}

%%%%%%%%%%%%%%%%%%%%%%%%%%%%%%%%%%%%%%%%%%%%%%%%%%%%
\z[{\begin{wrapfigure}[6]{r}{0.3\textwidth}%
		\vspace{-10pt}%
		\begin{flushright}%
 			\includegraphics[width=0.3\textwidth]{2.1.6}%
		\end{flushright}%
	\end{wrapfigure}%
	}]%
{На уроке физкультуры Петя и Маша бежали вместе по прямой дорожке, стартовав от школы. Затем Петя побежал быстрее, а Маша пошла. Через некоторое время дети одновременно повернули обратно и достигли школы так же одновременно. Графики зависимости проекции скорости детей на направление дорожки от времени даны на рисунке. Постройте график зависимости расстояния между Петей и Машей от времени.}
%
{$v = 50\,\text{м/мин}$}

%%%%%%%%%%%%%%%%%%%%%%%%%%%%%%%%%%%%%%%%%%%%%%%%%%%%
\z[{\begin{wrapfigure}[6]{r}{0.3\textwidth}%
		\vspace{-0pt}%
		\begin{flushright}%
 			\includegraphics[width=0.3\textwidth]{2.1.7}%
		\end{flushright}%
	\end{wrapfigure}%
	}]%
{Два мальчика, гуляя с собакой, вместе вошли в парк и побежали наперегонки по прямой дорожке к речке, до которой оставалось $1000\,\text{м}$ пути. Собака, не желая обижать никого из мальчиков, всё время бежала точно посередине между ними. Зависимости скоростей мальчиков от времени приведены на графике. К сожалению, единицы измерения по одной из осей не сохранились, но известно, что за $10\,\text{мин}$ мальчики добежали до речки. Определите минимальную скорость собаки за первые $8\,\text{мин}$ прогулки.}
%
{}